\subsection{LUDB and PTB Diagnostic reference-bounded comparisons}
\label{sec:ludb-ptb}
The completed LUDB study scored all 200 ECGs with no technical non-scoring. Its original joint policy accepted 152 and withheld 48. However, reference availability limits which endpoints can be evaluated. Principal-rhythm metadata gave 142 NSR references, 55 other recognized rhythms and three unknowns for the P8 NSR contrast. The NSR statement does not imply that the ECG is normal in every other respect. LVH had 108 positives and 92 unknowns; the unknowns are not negative controls. PAC had seven positives and 193 unknowns. Infarct-territory and ISC\_ performance was not estimable from these metadata.

On common endpoints with both positive and negative references, LUDB P6 AFIB/SB macro-F1 was 0.6949 on the full eligible source versus 0.8778 in the 59-ECG positive-profile subset. P8 AFIB/SB/NSR macro-F1 was 0.7075 on the full source versus 0.7407 in its 188-ECG restricted subset. These contrasts use matched endpoints and unchanged scores; they do not represent the original six- or eight-class headline metric. No P4 macro-F1 was calculated because none of its endpoints had adequate positive and explicit negative references in LUDB. The absence of a metric reflects a reference limitation rather than a zero score.

The separate PTB Diagnostic analysis retained all 549 ECGs from 290 subjects and scored them before joining reference information. Table~\ref{tab:ptb-transport} reports the subset with known references for each endpoint. Original P4 ASMI and IMI cutoffs were unchanged. ALMI and ILMI used previously fitted experimental PTB-XL heads and do not belong to the released P4 specialist. For the reported intervals, repeated ECGs were grouped by subject/exact-signal component. These source-specific retrospective analyses do not establish the independence of every historical training ancestor.

\begin{table}[tbp]
\caption{Completed retrospective PTB Diagnostic territory comparisons. All 549 ECGs were technically scored; other references remain unknown for each endpoint.\label{tab:ptb-transport}}
\centering\small
\begin{tabular}{@{}p{.20\linewidth}rrrp{.21\linewidth}p{.19\linewidth}@{}}
\toprule
Endpoint / score & Positive & Control & Unknown & TP/FN/FP/TN & F1 (cluster 95\% interval) \\
\midrule
ASMI / released P4 & 77 & 80 & 392 & 63/14/20/60 & 0.788 [0.672, 0.870] \\
IMI / released P4 & 89 & 80 & 380 & 66/23/21/59 & 0.750 [0.632, 0.833] \\
ALMI / experimental & 43 & 80 & 426 & 21/22/1/79 & 0.646 [0.523, 0.771] \\
ILMI / experimental & 53 & 80 & 416 & 22/31/1/79 & 0.579 [0.345, 0.764] \\
\bottomrule
\end{tabular}
\end{table}

The ECGs with explicit control references came from 52 source subjects. ASMI and IMI positives came from 27 and 30 subjects, respectively; ALMI and ILMI positives came from 16 and 22. These are not four independent cohorts because healthy controls and some subject information are shared. The compatible territory comparisons are narrower in scope than a comprehensive infarction detector. LVH and ISC\_ were not promoted to territory outcomes, and an unknown record was not converted to a disease-negative reference.

\subsection{Earlier and later morphology experiments}
\label{sec:morph-history}
The earlier fixed-representation paired-head study reported changes of $+0.72$, $+1.46$ and $+0.06$ percentage points for P6, P8 and P4, respectively. Its profile-specific 95\% intervals were $[-0.40,2.06]$, $[0.16,2.76]$ and $[-1.63,1.57]$ points. A later independent recalculation from the sealed aggregates reproduced the point estimates and bootstrap draws. In a separate post hoc within-group arm-swap analysis, Holm-adjusted $p$ values were 0.5681, 0.1356 and 0.9626 across the three profiles. Although P8's original unadjusted interval excluded zero, the post hoc family-wise test did not support a robust improvement. P6 and P4 remained inconclusive in that experiment; a nonsignificant comparison does not demonstrate equivalence.

The subsequent five-seed study used equal-capacity heads and fixed backbone representations, with the prespecified relevance margin described above. Its simultaneous intervals were within that margin for all three profiles. This supports a conclusion restricted to these fitted heads rather than a universal statement that morphology is ineffective. Seed-level differences were unstable for rare endpoints: the descriptive P6/P8 1AVB F1 differences ranged from $-5.38$ to $+4.60$ points, and P4 IMI from $-5.33$ to $+6.44$. Neither range is an uncertainty interval for clinical benefit.

The reconstructed runtime was also sensitive to preprocessing. On the 16-ECG instrumented batch, removing the bandpass changed maximum P4 probabilities by 0.5157, omitting standardization by 0.8491 and dropping lead II by 0.3225. CPU versus the recorded GPU calculation changed maximum P4 probability by 0.000531, while repeated fixed-batch GPU forwards agreed exactly. These perturbation maxima describe a small technical probe rather than comparative diagnostic performance. Repeatability of one runtime therefore cannot resolve the unavailable historical loader identity.
